\documentclass[10pt,a4paper]{amsart}
\usepackage[margin=19mm]{geometry}
\usepackage{amsmath,amssymb,mathtools}
\usepackage[scr=rsfs,cal=euler]{mathalfa}
\usepackage{xfrac}
\usepackage[unicode,hidelinks,bookmarks=false]{hyperref}
\newcommand{\ie}{i.e.\ }
\newcommand{\cf}{cf.\ }
\newcommand{\resp}{respectively\ }
\newcommand{\Subsection}{\subsection}
\providecommand{\nobreakdash}{\nobreak\hbox{-}}
\setlength{\emergencystretch}{2em}
\multlinegap0pt
\usepackage{amssymb}
\usepackage{enumerate}
\usepackage{xcolor}
\newtheorem{corollary}{Corollary}
\usepackage{cleveref}
\crefname{corollary}{Corollary}{Corollarys}
\renewcommand{\AA}{\mathbb{A}}
\newcommand{\Pic}{\operatorname{Pic}}
\title{Relative $\AA^{1}$-contractibility of smooth schemes}
\author{Adrien Dubouloz, Krishna Kumar Madhavan Vijayalakshmi, Paul Arne {\O}stv{\ae}r}
\date{Source: arXiv:2605.07638v1 (CC BY 4.0)}
\begin{document}
\maketitle

\noindent\emph{Source and licence.} Relative $\AA^{1}$-contractibility of smooth schemes. Version arXiv:2605.07638v1 (CC BY 4.0), distributed by arXiv under the Creative Commons Attribution 4.0 International licence, http://creativecommons.org/licenses/by/4.0/, as recorded in the arXiv licence metadata for this exact version and shown on the abstract page https://arxiv.org/abs/2605.07638v1. Full licence notice: \texttt{licenses/arxiv-2605.07638-CC-BY-4.0.txt}.\par\smallskip

\noindent\textit{Editorial scope.} Counted unit: the article's corollary corollary (source0.tex lines 1054--1060). The article's complete original proof of it is delivered with it (source0.tex lines 1062--1084, one \texttt{proof} environment of 23 lines). No mathematics is written, repaired or re-derived: the statement and the proof are the article's own lines, in the article's own notation, and no retained line is substituted or re-flowed.  Retained uncounted as the source-local context the proof needs: lines 1040--1042.  Nothing was omitted from any retained span, so the package carries no omission record.  No retained line contains a citation, so the wrapper carries no bibliography and no bibliography entry.  No retained line contains a figure environment, an image-inclusion command or drawing code, so no image is delivered and the package declares no asset dependency.  Editorial formatting only, in addition: an AMS \texttt{amsart} preamble that restates the article's own declarations and macros used by the retained lines, namely the environments \texttt{corollary} on separate counters, together with the standard packages those lines need.  The retained environments print in this document as Corollary 1 in order of appearance, whereas the article prints the same items with the numbering of its own full document.  The complete native source of the article as distributed in the arXiv e-print for this exact version is bundled unchanged as \texttt{sources/200101/source0.tex}.
\begin{corollary}\label{A1unique:DD}
Let $S$ be a noetherian normal scheme. A smooth morphism $f:X\to S$ of finite presentation and relative dimension $1$ is isomorphic to $\AA^1_S$ if and only if $f$ is an $\AA^1$-weak equivalence in $\mathrm{Spc}_S^{\AA^1}$, the relative canonical sheaf $\omega_f=\Omega_{X/S}$ is trivial, and $f$ admits a section.
\end{corollary}

\begin{corollary}
Let $f:X\to S$ be a smooth morphism of finite presentation and relative dimension $1$ over a normal scheme $S$. If there exists $n\geq 1$ such that
\[
X\times_S \AA^n_S \cong \AA^{n+1}_S,
\]
then $X$ is $S$-isomorphic to $\AA^1_S$.
\end{corollary}

\begin{proof}
From the given isomorphism we obtain
\[
\mathrm{pr}_1^*\omega_f
\;\cong\;
\omega_{X\times_S \AA^n_S/S}
\;\cong\;
\omega_{\AA^{n+1}_S/S}
\;\cong\;
\mathcal{O}_{X\times_S \AA^n_S},
\]
where $\mathrm{pr}_1$ denotes the projection. Since $S$ is normal and $f$ is smooth, $X$ is normal, and the pullback homomorphism
\[
\mathrm{pr}_1^*:\Pic(X)\to \Pic(X\times_S \AA^n_S)
\]
is an isomorphism. It follows that $\omega_f$ is trivial. Moreover, $f$ admits a section obtained by composing the zero section $s_0:S\to \AA^{n+1}_S$ with the projection $\AA^{n+1}_S\cong X\times_S \AA^n_S\to X$. Finally, since both projections
\[
X\times_S \AA^n_S\to X
\quad\text{and}\quad
\AA^{n+1}_S\to S
\]
are $\AA^1$-weak equivalences, it follows that $f:X\to S$ is an $\AA^1$-weak equivalence. The conclusion now follows from \Cref{A1unique:DD}.
\end{proof}

\end{document}
